\subsection{Sorted modular pair streams}
\label{sec:streams}

A stream will enumerate a Cartesian product while retaining only its
smaller input arrays. An item consists of its actual, possibly signed,
weight and a constant number of index masks. Write $\omega(a)$ for
the actual integer weight of an item $a$. Repeated occurrences
remain separate even if all their fields agree. Each item occupies
constantly many words; array indices have $O(n)$ bits. We assume that
weights, moduli, keys, and indices fit in the stated constant-word
representation. Section~\ref{sec:implementation} verifies these
conditions for all uses below.

This is the sorted pair-sum method associated with Schroeppel and
Shamir~\cite{SS81}; see also the explicit four-list description in
Nederlof and W\k{e}grzycki~\cite[full version, Lemmas~3.3 and~A.1]{NW21}.
We give the modular ordering and stopped-execution argument locally,
since those details determine the cost of regenerating a dictionary.
A generator's mutable state identifies its next occurrence and can be
copied while sharing its immutable input arrays.

\begin{lemma}[Filtered streams and tie joins]
\label{lem:streams}
Let $A_1,A_2$ be stored constant-word item arrays whose indices have
$O(n)$ bits. Let $p\mid P$ be positive integers, and let $c,v$ be
specified integer residues. Assume that the data occupy constantly
many words each and that arithmetic, comparisons, and indexed access
on them take $O(1)$ word operations in the stated model. The pairs $(a,b)\in A_1\times A_2$
whose actual weights sum to $c\pmod p$ can be streamed in
nondecreasing order of either
\[
 (\omega(a)+\omega(b))\bmod P
 \quad\text{or}\quad
 (v-\omega(a)-\omega(b))\bmod P,
\]
where all residues are normalized to $\{0,\ldots,P-1\}$.
Initialization costs $O(n(|A_1|+|A_2|+1))$ word operations, each
subsequent advance costs $O(n)$, and the state together with the
input arrays uses $O(|A_1|+|A_2|+1)$ words.

Consider two such streams of lengths $N_1,N_2$, with combined states
and arrays of at most $S$ words. If their equal-key Cartesian
products contain $Z$ pairs in all, they can report all these pairs
using $O(S+B)$ words and, apart from initialization and the work
performed on reported pairs,
\begin{equation}
\label{eq:tie-cost}
 O\left(B+n(N_1+N_2+Z)+\frac SB Z\right)
\end{equation}
operations. Here $B\ge1$ is an integer buffer capacity. In
particular, if $S\le\poly(n)B$, the bound is
$\poly(n)(B+N_1+N_2+Z)$. The same upper bounds hold for an execution
stopped partway, with $Z$ still counting \emph{all potential} matches.
No coprimality assumption is needed.
\end{lemma}

\begin{proof}
First sort $A_2$ lexicographically by its residue modulo $p$ and
then its residue modulo $P$. For a fixed $a\in A_1$, binary
search finds the contiguous slice with residue
$c-\omega(a)\pmod p$. This slice contains exactly the eligible
partners of $a$, and is ordered by their $P$-residues, including
repetitions. Omit the row if the slice is empty.

For a direct key, put $r=\omega(a)\bmod P$. Start at the first
entry of the slice whose residue is at least $P-r$; if there is
none, start at the beginning. Traverse forward cyclically for exactly
the slice length. The entries at or above the pivot give the wrapped
values $r+(\omega(b)\bmod P)-P$ first, followed by the unwrapped
values $r+(\omega(b)\bmod P)$.
Both portions are nondecreasing and the former values are smaller
than the latter. This includes $r=0$, when the pivot search simply
returns the beginning.

For a complementary key put $K=(v-\omega(a))\bmod P$.
Start at the last entry whose residue is at most $K$; if there is
none, start at the last entry of the slice. Traverse in reverse
cyclic order for exactly one slice length. Residues at most $K$
produce the values $K-(\omega(b)\bmod P)$ in nondecreasing
order. The residues above $K$ follow and produce the larger values
$K-(\omega(b)\bmod P)+P$, also in nondecreasing order.
Choosing the first pivot in the direct case and the last pivot in
the complementary case preserves every repeated residue.

Merge the rows with a min-heap having at most one current pair per
item of $A_1$. A row stores its slice endpoints, current position,
and remaining length; a heap entry has constantly many words.
Comparison sorting, binary searches, and heap operations have depth
$O(n)$ because array indices have $O(n)$ bits. This proves the
initialization, advance, and state bounds. The construction permits
repeated prime factors and other common factors. Empty input arrays
give an empty stream.

To join two streams, advance past unequal keys until an equal key is
found. For that key, read up to $B$ second-stream items into a reusable
buffer. If this exhausts the second tie group, pair its buffered items
with each consecutive first-stream item of the key. No generator state
is copied in this case. Reuse the buffer by its current length, without
clearing $B$ positions for a short group.

If the second tie group is longer, save the mutable second-generator
state immediately after the first $B$ items. The snapshot includes
its current lookahead item, heap, row cursors, and control state, but
shares its immutable arrays. For each first-stream item of the key,
report its pairs with the buffer, restore the saved state when needed,
and replay the remaining suffix. After the final replay the second
stream is already beyond the tie group.

A save or restore costs $O(S)$. It is made only for a first item
having at least $B$ potential partners; at most a constant number of
such operations per first item is required. Their total cost is
therefore $O(SZ/B)$. Replayed advances are charged to matching pairs;
other advances are charged to the original stream lengths. The
$O(n)$ cost per advance and the optional $O(B)$ initial buffer
allocation give~\eqref{eq:tie-cost}. If processing stops immediately
after a snapshot or restore, the same charge remains valid because
$Z$ counts potential pairs, including pairs not reached. One snapshot,
the buffer, and the two generators give the space bound.
\end{proof}

An actual-weight stream needs no modular pivots. Sort the second array
by actual weight, add each first-array weight to that sorted row, and
heap-merge the rows. Reading this ordered stream while remembering
only the last weight counts its distinct values without storing the
Cartesian product. This will let the representation trial measure a
block's deficiency exactly from two small half-block lists.

\subsection{Bounded dictionaries}

Streaming saves space only if we also control what is retained.
For deterministic deduplication, use a binary trie on a fixed-width
signed actual weight, followed by the relevant mask bits when a mask
is part of the key. Give zero its canonical nonnegative sign. In our
applications the complete key has $O(u)$ bits. Searching or inserting
it takes $O(u)$ operations and introduces at most $O(u)$ nodes,
with constantly many words per node.

Keep a separate append-only array of the distinct entries if subsequent
array access is needed. A duplicate changes neither the array length
nor its mathematical content. On finding a new key that would increase
the array beyond capacity $C$, report overflow. Even if nodes for this
first rejected key have been allocated, the trie has at most
$O(u(C+1))$ nodes. All nodes, addresses, and appended records count as
workspace. A weight-only dictionary is appropriate for disjoint
compression blocks; a weight--mask dictionary will be required when
compatibility depends on shared indices.

For the concrete implementations below, raw list generation can be
done in $O(n^2)$ operations per item, sorting and heap work in $O(n)$
per item or advance, and dictionary operations in $O(u)$ per key.
Whenever the combined stream storage is $O(B)$, the term depending
on $Z$ in~\eqref{eq:tie-cost} is $O(nZ)$, before the separately
charged processing of those matches. The buffer and stream-length
terms remain separate. These explicit
degrees will be used when the task costs are bounded. The argument
does not assign a cost merely from the number of distinct entries
eventually retained.
